\chapter{GIỚI THIỆU}

\section{Đặt vấn đề}

Flappy Bird có luật chơi đơn giản nhưng đặt ra nhiều yêu cầu điển hình của một hệ thống nhúng đa phương tiện: cập nhật vật lý theo thời gian cố định, phản hồi input nhanh, vẽ lại toàn màn hình, tránh xé hình, phát âm thanh liên tục và sử dụng bộ nhớ hợp lý. Trên STM32F429, các yêu cầu này phải được giải quyết trong giới hạn \SI{2}{\mega\byte} Flash, \SI{256}{\kilo\byte} SRAM nội và không có hệ điều hành.

Màn hình $240\times320$ ở định dạng ARGB8888 cần $240\times320\times4=307\,200$ byte cho một khung hình. Chỉ một framebuffer đã vượt quá dung lượng của một bank SRAM chính, còn hai framebuffer cần $614\,400$ byte. Vì vậy, việc khai thác SDRAM \SI{8}{\mega\byte} gắn trên board là điều kiện cần để thực hiện double buffering.

Âm thanh cũng yêu cầu luồng dữ liệu đều đặn. Nếu CPU tự ghi từng mẫu vào I2S, thời gian dành cho game và đồ họa sẽ bị phân mảnh. DMA circular giải quyết vấn đề truyền liên tục, nhưng phần mềm vẫn phải nạp kịp hai nửa buffer và tránh làm quá nhiều việc trong ISR.

\section{Mục tiêu đề tài}

Đề tài hướng tới các mục tiêu sau:

\begin{itemize}
  \item Xây dựng game Flappy Bird chạy độc lập trên STM32F429I-Discovery, không cần máy tính khi vận hành.
  \item Duy trì nhịp cập nhật ổn định khoảng \SI{50}{FPS} bằng timer phần cứng.
  \item Hiển thị đồ họa mượt trên LCD ILI9341, sử dụng SDRAM và double buffering.
  \item Nhận thao tác từ màn hình cảm ứng và nút nhấn vật lý có chống dội.
  \item Phát nhạc nền và hiệu ứng đồng thời qua MAX98357 bằng I2S và DMA.
  \item Tổ chức mã nguồn thành các module phần cứng, game, âm thanh và input có giao diện rõ ràng.
\end{itemize}

\section{Nền tảng phần cứng và công cụ}

\begin{table}[H]
  \centering
  \caption{Nền tảng sử dụng trong project}
  \label{tab:platform}
  \begin{tabularx}{\textwidth}{>{\bfseries}p{4cm}X}
    \toprule
    Thành phần & Mô tả \\
    \midrule
    Vi điều khiển & STM32F429ZIT6, ARM Cortex-M4F, xung tối đa \SI{180}{\mega\hertz}, Flash \SI{2}{\mega\byte}. \\
    Kit phát triển & STM32F429I-Discovery với LCD TFT 2.4 inch, SDRAM \SI{8}{\mega\byte} và cảm ứng STMPE811. \\
    Âm thanh & MAX98357, nhận dữ liệu I2S và khuếch đại class-D cho loa ngoài. \\
    IDE & STM32CubeIDE, GNU Tools for STM32, ngôn ngữ C11. \\
    Cấu hình & STM32CubeMX thông qua file \code{FlappyBird.ioc}. \\
    Thư viện & STM32 HAL, CMSIS và BSP STM32F429I-Discovery. \\
    \bottomrule
  \end{tabularx}
\end{table}

\section{Phạm vi và mã nguồn}

Mã nguồn của project được tổ chức thành các module ứng dụng chính gồm \code{main.c}, \code{game.c}, \code{audio.c}, driver LCD/SDRAM, driver cảm ứng và STMPE811. Mã nguồn quản lý nhóm được lưu tại \href{\RepositoryUrl}{GitHub của project}.

Project không sử dụng TouchGFX hoặc RTOS. Giao diện được vẽ trực tiếp bằng BSP LCD, còn lịch trình được điều phối bằng superloop, TIM6, SysTick và các cờ ngắt. Báo cáo không đánh giá chất lượng âm thanh bằng thiết bị chuyên dụng và không đưa ra số liệu FPS đo bằng oscilloscope khi chưa có bằng chứng thực nghiệm.

\section{Bố cục báo cáo}

Sau phần giới thiệu, Chương 2 trình bày phân công công việc. Chương 3 phân tích phần cứng và các ngoại vi. Chương 4 mô tả kiến trúc phần mềm, thuật toán game và âm thanh. Chương 5 tổng hợp kết quả build, tài nguyên, khó khăn và đánh giá. Chương 6 nêu kết luận cùng hướng phát triển.
